% ---------------------------------------------------------------
% DRAFT OUTLINE. The content below is rendered as bullet lists so that
% it compiles and can be read/commented by co-authors. Replace each
% itemize block with prose as the section is written.
%
% Class option note: [preprint] gives a readable single-column draft.
% Switch to [5p] for the two-column submission format.
% ---------------------------------------------------------------
\documentclass[preprint,11pt]{elsarticle}
% \documentclass[5p]{elsarticle}   % <- submission format

\usepackage{amsmath}
\usepackage{amssymb}
\usepackage{graphicx}
\usepackage{xcolor}
\usepackage{todonotes}
\usepackage{overpic}
\usepackage{wasysym}
\usepackage[dvipsnames]{xcolor}
\usepackage{enumitem}
\usepackage[colorlinks = true, linkcolor = blue, urlcolor  = blue, citecolor = blue]{hyperref}

% Compact nested lists for the outline
\setlist[itemize]{noitemsep, topsep=3pt, parsep=2pt, leftmargin=1.4em}

% ---------------------------------------------------------------
% Macros
% ---------------------------------------------------------------
\newcommand{\FESTIM}{\textsc{festim}}
\newcommand{\TMAP}{\textsc{tmap}}
\newcommand{\MHIMS}{\textsc{mhims}}
\newcommand{\HYPERION}{\textsc{hyperion}}
\newcommand{\Da}{\mathrm{Da}}
\newcommand{\kB}{k_\mathrm{B}}
\newcommand{\half}{\tfrac{1}{2}}
\newcommand{\flibe}{\mathrm{FLiBe}}
\newcommand{\cm}[1]{c^{\mathrm{m}}_{#1}}   % metal-side (atomic) concentration
\newcommand{\cs}[1]{c^{\mathrm{s}}_{#1}}   % salt-side (molecular) concentration
\newcommand{\Bra}{\mathcal{B}}             % branching ratio

% Draft annotation macros -- delete these (and their uses) before submission
\newcommand{\TODO}[1]{\textcolor{BrickRed}{\textbf{[TODO]}~#1}}
\newcommand{\VERIFY}[1]{\textcolor{RedOrange}{\textbf{[VERIFY]}~#1}}
\newcommand{\KEY}[1]{\textcolor{RoyalBlue}{\textbf{[KEY POINT]}~#1}}

% ---------------------------------------------------------------
% TODO(journal): elsarticle + \journal{Nuclear Fusion} was inconsistent
% (NF is IOP). With the molten-salt / fuel-cycle framing the natural
% Elsevier targets are:
%   - Fusion Engineering and Design      <- set below
%   - International Journal of Hydrogen Energy (where FESTIM v2 landed)
%   - Nuclear Materials and Energy
% ---------------------------------------------------------------
\journal{Fusion Engineering and Design}
% ---------------------------------------------------------------

\begin{document}

\begin{frontmatter}

% TODO(title): pick one
%  (A) Beyond local thermodynamic equilibrium: kinetic interface models
%      for multi-isotope hydrogen transport in metal--molten-salt systems
%  (B) When chemical potential continuity fails: kinetic interface models
%      for hydrogen isotope transport
%  (C) as below
\title{Kinetic interface models for hydrogen isotope transport:\\
       local thermodynamic equilibrium as a limiting case}

\author[mit]{Remi Delaporte-Mathurin}
\author[mit]{James Dark}
% TODO(authors): add HYPERION experimental team + salt chemistry co-authors

\address[mit]{Massachusetts Institute of Technology, Cambridge, MA, USA}

\begin{abstract}
\TODO{write the abstract last. Points to hit:}
\begin{itemize}
    \item Macroscopic hydrogen transport codes (\FESTIM, \TMAP, \MHIMS) impose
          local thermodynamic equilibrium (LTE) at interfaces: continuity of
          chemical potential, realised as a per-species Dirichlet constraint.
    \item We formulate a general kinetic interface framework: a set of
          reversible interfacial reaction channels obeying mass action, with
          detailed balance fixing the rate ratios.
    \item LTE is recovered analytically as the fast-kinetics limit of a
          \emph{single} channel --- in both the Sieverts/Sieverts and the
          Sieverts/Henry forms.
    \item A Damk\"ohler criterion identifies when kinetics matter \emph{within}
          a channel; a branching ratio identifies when no LTE form exists at all.
    \item Applied to \HYPERION{} (Ni / molten FLiBe / gas): hydrogen leaving the
          Ni can recombine to $\mathrm{H_2}$ or be fluorinated to HF, giving a
          permeating flux that is a kinetically-set mixture of two carriers. The
          apparent interfacial law is neither Sieverts nor Henry, but drifts
          between them with loading, temperature and salt redox state.
    \item With H and T present the salt carries five species
          ($\mathrm{H_2}$, HT, $\mathrm{T_2}$, HF, TF) fed by two metal-side
          species through a bilinear map: a per-species LTE condition is not
          merely inaccurate, it is ill-posed. Isotope swamping and interfacial
          fractionation are predicted; neither is expressible under LTE.
    \item Flux speciation becomes a direct simulation output, comparable with
          \HYPERION{} bubbler and mass-spectrometry diagnostics.
\end{itemize}
\end{abstract}

\begin{keyword}
    hydrogen isotope transport \sep
    material interfaces \sep
    local thermodynamic equilibrium \sep
    \FESTIM \sep
    molten salt \sep
    FLiBe \sep
    nickel \sep
    isotope exchange \sep
    kinetic model
\end{keyword}

\end{frontmatter}

% ===============================================================
\section{Introduction}
\label{sec:intro}
% ===============================================================

\begin{itemize}
    \item Hydrogen isotope transport controls fuel retention, permeation losses
          and tritium accountancy in fusion systems. In liquid-breeder concepts
          the tritium path crosses at least one metal/liquid interface before
          reaching an extraction or release point.

    \item Macroscopic transport codes (\FESTIM{} \cite{dark_festim_2026},
          TMAP8 \cite{simon_moose-based_2025},
          \MHIMS{} \cite{hodille_macroscopic_2015}) close such interfaces with
          local thermodynamic equilibrium (LTE), i.e.\ continuity of the
          chemical potential of the transported species. In practice this is
          imposed as an algebraic (Dirichlet-type) constraint relating the two
          interfacial concentrations:
          \begin{itemize}
              \item Sieverts/Sieverts: $c_A/K_{S,A} = c_B/K_{S,B}$
              \item Sieverts/Henry: $c_B = K_{H,B}\,(c_A/K_{S,A})^2$
          \end{itemize}

    \item Three assumptions are hidden in that closure, and they are rarely
          stated together:
          \begin{itemize}
              \item[(H1)] interfacial equilibration is fast compared with bulk
                    transport;
              \item[(H2)] a \emph{single} exchange pathway connects the two sides;
              \item[(H3)] the identity of the carrier species on each side is
                    known a priori and fixed, so that a single solubility law
                    applies.
          \end{itemize}

    \item (H1) is the assumption usually scrutinised, and it is known to fail:
          \begin{itemize}
              \item Hodille et al.\ \cite{hodille_molecular_2022}: MD at the
                    Be/BeO interface shows LTE is not reached even at 1500~K.
              \item Silva-Sol\'is et al.\ \cite{silva-solis_hydrogen_2026}:
                    DFT+MRE for W/Cu shows the steady state is not the
                    thermodynamic equilibrium state whenever a net flux is
                    carried.
          \end{itemize}

    \item \KEY{(H2) and (H3) have received almost no attention, and this paper
          argues they are the more consequential failures.} The motivating
          system is a metal in contact with a fluoride melt.

    \item Metal/molten-salt interfaces: hydrogen leaving the metal lattice has
          more than one chemical fate. It may recombine to $\mathrm{H_2}$ and
          dissolve physically (Henry's law, low solubility, high mobility), or
          it may be oxidised by the melt to HF and dissolve chemically (much
          higher solubility, corrosive, redox dependent). The branching between
          the two is kinetic and depends on the salt redox state.

    \item This ambiguity is not hypothetical. Calderoni et al.\
          \cite{calderoni_measurement_2008}, measuring T permeation through a Ni
          membrane into FLiBe, concluded that tritium absorbed atomically in the
          nickel does \emph{not} simply recombine at the Ni/FLiBe interface, and
          could not distinguish between transport as T bound to
          $\mathrm{BeF_4^-}$, as HT, or both. Nearly two decades later the
          speciation at that interface is still an open question, and reported
          FLiBe transport properties for $\mathrm{H_2}$/HF and their
          isotopologues still span orders of magnitude \cite{libra_2022}.

    \item The ambiguity has already contaminated the data record. The solubility
          constant reported in \cite{calderoni_measurement_2008} carries units
          inconsistent with the dissolution law it is routinely used with
          downstream (Sec.~\ref{sec:units}). This is not a typographical slip:
          \emph{a solubility constant has no law-independent units}.
          Sieverts-type dissolution gives
          $\mathrm{mol\,m^{-3}\,Pa^{-1/2}}$, Henry-type gives
          $\mathrm{mol\,m^{-3}\,Pa^{-1}}$, and the two cannot be interconverted
          by any fixed factor. Reporting a number therefore presupposes the
          speciation --- exactly the quantity in question. We treat this as a
          motivating symptom rather than a criticism of the measurement: it is
          what assumption (H3) looks like when it fails in practice.

    \item It gets structurally worse with more than one isotope. Two mobile
          atomic species in the metal (H, T) feed five carriers in the salt
          ($\mathrm{H_2}$, HT, $\mathrm{T_2}$, HF, TF). The map is quadratic and
          non-diagonal: $c_\mathrm{HT}$ depends on the \emph{product} of the two
          metal-side concentrations. No per-species chemical-potential
          continuity condition can express this --- LTE here is not inaccurate,
          it is ill-posed.

    \item Gap: no macroscopic hydrogen transport code implements a general
          kinetic interface condition with competing reaction channels and
          multiple isotopologues.

    \item This work:
          \begin{enumerate}[label=(\roman*)]
              \item formulates a general kinetic interface framework based on
                    reversible reaction channels and mass action, with detailed
                    balance fixing rate ratios from thermodynamics;
              \item implements it in \FESTIM{} \cite{dark_festim_2026};
              \item recovers LTE analytically as the fast-kinetics
                    single-channel limit, in both Sieverts/Sieverts and
                    Sieverts/Henry forms;
              \item derives a Damk\"ohler criterion (within a channel) and a
                    branching ratio (between channels) delimiting LTE validity;
              \item applies the framework to the \HYPERION{} Ni/FLiBe/gas
                    experiment \cite{hyperion_facility}, in single-isotope and
                    H/T modes.
          \end{enumerate}

    \item \textbf{Scope note.} Interfaces are treated here as reactive surfaces
          without their own stored inventory. Intermediate trapped states
          localised on the interface --- and the codimension-1 formulation
          needed to carry an interfacial ODE alongside the bulk PDEs --- are the
          subject of a companion paper \cite{festim_codim1_inprep}. The Be/BeO
          and W/Cu interface-trapping cases of
          Refs.~\cite{hodille_molecular_2022,silva-solis_hydrogen_2026} belong
          there.
          \TODO{write this as two explicit sentences in the final prose so
          referees do not ask.}

    \item \TODO{paper outline paragraph.}
\end{itemize}

% ===============================================================
\section{Interface conditions for hydrogen transport}
\label{sec:models}
% ===============================================================

% ---------------------------------------------------------------
\subsection{Local thermodynamic equilibrium and its assumptions}
\label{sec:lte}
% ---------------------------------------------------------------

\begin{itemize}
    \item Two subdomains $\Omega_A$, $\Omega_B$ sharing an interface $\Gamma$.
          Bulk transport $\partial_t c = \nabla\!\cdot\!(D\nabla c) + S$ follows
          \cite{dark_festim_2026}; traps are omitted here for clarity.

    \item LTE is the equality of the chemical potential of dissolved hydrogen on
          both sides of $\Gamma$. Written in terms of concentrations it takes a
          different algebraic form depending on the dissolution law:
\end{itemize}

\begin{equation}
    \frac{c_A\big|_\Gamma}{K_{S,A}(T)}
    = \frac{c_B\big|_\Gamma}{K_{S,B}(T)}
    \qquad \text{(Sieverts / Sieverts)}
    \label{eq:lte_ss}
\end{equation}

\begin{equation}
    c_B\big|_\Gamma
    = K_{H,B}(T)\left(\frac{c_A\big|_\Gamma}{K_{S,A}(T)}\right)^{2}
    \qquad \text{(Sieverts / Henry)}
    \label{eq:lte_sh}
\end{equation}

\begin{itemize}
    \item Both are \emph{algebraic} constraints: they fix a ratio (or a power
          law) between interfacial concentrations, independently of the flux the
          interface is carrying. Consequences:
          \begin{itemize}
              \item no interfacial timescale: the interface responds
                    instantaneously;
              \item no interfacial resistance: $\Gamma$ never limits the flux;
              \item the exponent (1 or 2) must be chosen a priori, i.e.\ the
                    carrier species must be known before the simulation is run.
          \end{itemize}

    \item Formal point, following \cite{silva-solis_hydrogen_2026} Sec.~IV.C:
          thermodynamic equilibrium requires every microscopic forward rate to
          be balanced by its reverse (detailed balance), hence \emph{zero} net
          flux across $\Gamma$. A permeation experiment carries a nonzero net
          flux by construction. LTE is therefore always an approximation; the
          question is only how large the error is.

    \item Distinguish clearly, and keep this distinction throughout:
          \begin{itemize}
              \item steady state: $\partial_t c = 0$,
                    $\Gamma_\mathrm{in} = \Gamma_\mathrm{out} \neq 0$;
              \item equilibrium: $\partial_t c = 0$,
                    $\Gamma_\mathrm{in} = \Gamma_\mathrm{out} = 0$.
          \end{itemize}
          Equilibrium is a steady state; the converse is false.
\end{itemize}

% ---------------------------------------------------------------
\subsection{A general kinetic interface framework}
\label{sec:general}
% ---------------------------------------------------------------

\begin{itemize}
    \item Replace the algebraic constraint by a set of reversible reaction
          channels living on $\Gamma$. Let $X_i$ denote mobile atomic species on
          side $A$ ($i \in \{\mathrm{H,D,T}\}$) and $Y_\alpha$ the carriers on
          side $B$ ($\alpha \in \{\mathrm{H_2, HT, T_2, HF, TF},\dots\}$).
          Channel $r$ reads
          \begin{equation*}
              \sum_i \nu_{ir} X_i + \sum_\beta \lambda_{\beta r} Z_\beta
              \;\rightleftharpoons\; \sum_\alpha \mu_{\alpha r} Y_\alpha
          \end{equation*}
          where $Z_\beta$ are non-hydrogenic salt constituents entering the
          stoichiometry ($\mathrm{F^-}$, oxidant/reductant of the redox buffer).
\end{itemize}

\begin{equation}
    w_r = k_r^{+}(T)\prod_i \left(\cm{i}\big|_\Gamma\right)^{\nu_{ir}}
                     \prod_\beta a_\beta^{\lambda_{\beta r}}
        - k_r^{-}(T)\prod_\alpha \left(\cs{\alpha}\big|_\Gamma\right)^{\mu_{\alpha r}}
    \label{eq:mass_action}
\end{equation}

\begin{itemize}
    \item Interface conditions: the net atomic flux of isotope $i$ out of $A$,
          and the production of carrier $\alpha$ into $B$, are sums over
          channels.
\end{itemize}

\begin{align}
    -D^{\mathrm{m}}_i \nabla \cm{i}\cdot \mathbf{n}\big|_\Gamma
        &= \sum_r \nu_{ir}\, w_r
    \label{eq:bc_metal}\\
     D^{\mathrm{s}}_\alpha \nabla \cs{\alpha}\cdot \mathbf{n}\big|_\Gamma
        &= \sum_r \mu_{\alpha r}\, w_r
    \label{eq:bc_salt}
\end{align}

\begin{itemize}
    \item Atom conservation across $\Gamma$ is enforced by the stoichiometry.
          \TODO{state the balance condition explicitly (sum over channels of
          $\nu$ versus sum of $\mu$ weighted by atoms per carrier).}

    \item \textbf{Detailed balance.} Each channel independently satisfies
          $k_r^{+}/k_r^{-} = K_r(T)$, with $K_r$ the equilibrium constant of
          that channel obtained from the solubilities and the relevant formation
          free energies. Thermodynamics fixes the \emph{ratio}; atomistics or
          experiment fixes the \emph{magnitude}. This is what makes the
          framework a strict generalisation of LTE rather than a competing
          model.

    \item The framework is agnostic to what is on each side: the gas/solid
          recombination--dissociation boundary condition already in \FESTIM{} is
          the special case $A = \mathrm{gas}$.
\end{itemize}

% ---------------------------------------------------------------
\subsection{Model 1 --- first-order exchange}
\label{sec:model1}
% ---------------------------------------------------------------

\begin{itemize}
    \item Single channel, $X_A \rightleftharpoons X_B$: atomic transfer across a
          metal/metal interface. Robin-type flux condition:
\end{itemize}

\begin{equation}
    \phi = k^{+} c_A\big|_\Gamma - k^{-} c_B\big|_\Gamma
    \label{eq:model1_flux}
\end{equation}

\begin{itemize}
    \item Mass conservation:
          $-D_A \,\partial_n c_A = D_B\, \partial_n c_B = \phi$.
    \item Detailed balance ($\phi = 0 \Rightarrow$ Eq.~\ref{eq:lte_ss}):
\end{itemize}

\begin{equation}
    \frac{k^{+}}{k^{-}} = \frac{K_{S,B}(T)}{K_{S,A}(T)}
    \label{eq:detailed_balance_1}
\end{equation}

\begin{itemize}
    \item LTE limit: $k^{+}, k^{-} \rightarrow \infty$ at fixed ratio recovers
          Eq.~(\ref{eq:lte_ss}) exactly. Derivation in
          App.~\ref{app:limits}; the finite-$k$ correction is
          $\mathcal{O}(\phi/k^{+}c)$ and is quantified in
          Sec.~\ref{sec:damkohler}.
    \item Units: $k^{+}, k^{-}$ in $\mathrm{m\,s^{-1}}$.
\end{itemize}

% ---------------------------------------------------------------
\subsection{Model 2 --- recombination into a molecular carrier}
\label{sec:model2}
% ---------------------------------------------------------------

\begin{itemize}
    \item Single channel,
          $2\,X_A \rightleftharpoons Y_{A_2}$: atomic hydrogen on the metal
          side, molecular hydrogen dissolved on the liquid side. This is the
          metal/molten-salt analogue of surface recombination.
\end{itemize}

\begin{equation}
    w_{\mathrm{rec}} = k_{\mathrm{r}}^{+}\left(\cm{\mathrm{H}}\big|_\Gamma\right)^{2}
                     - k_{\mathrm{r}}^{-}\,\cs{\mathrm{H_2}}\big|_\Gamma
    \label{eq:model2_rate}
\end{equation}

\begin{itemize}
    \item Interface conditions:
          $-D^{\mathrm{m}}\partial_n \cm{\mathrm{H}} = 2 w_{\mathrm{rec}}$ and
          $D^{\mathrm{s}}\partial_n \cs{\mathrm{H_2}} = w_{\mathrm{rec}}$.
    \item Detailed balance. Setting $w_{\mathrm{rec}} = 0$ and using Sieverts in
          the metal ($c^{\mathrm{m}} = K_S\sqrt{P}$) with Henry in the salt
          ($c^{\mathrm{s}} = K_H P$):
\end{itemize}

\begin{equation}
    \frac{k_{\mathrm{r}}^{+}}{k_{\mathrm{r}}^{-}}
    = \frac{K_{H}(T)}{K_{S}(T)^{2}}
    \label{eq:detailed_balance_2}
\end{equation}

\begin{itemize}
    \item \KEY{Substituting Eq.~(\ref{eq:detailed_balance_2}) into
          $w_{\mathrm{rec}} = 0$ returns exactly the Sieverts/Henry LTE
          condition, Eq.~(\ref{eq:lte_sh}).} So the mixed-law interface
          currently available in \FESTIM{} v2 \cite{dark_festim_2026} is the
          fast-kinetics limit of Eq.~(\ref{eq:model2_rate}). Worth stating
          plainly: the two interface laws of Sec.~\ref{sec:lte} are not two
          separate physical models, they are the fast limits of two different
          channels.
\end{itemize}

% ---------------------------------------------------------------
\subsection{Model 3 --- competing channels}
\label{sec:model3}
% ---------------------------------------------------------------

\begin{itemize}
    \item Now let \emph{two} channels operate on the same interface. In a
          fluoride melt, hydrogen leaving the metal may either recombine
          (channel R, Eq.~\ref{eq:model2_rate}) or be oxidised by the melt
          (channel F):
          \begin{itemize}
              \item R:\quad $2\,\mathrm{H(m)} \rightleftharpoons \mathrm{H_2(s)}$
              \item F:\quad $\mathrm{H(m)} + \mathrm{F^-(s)} + h^{+}
                    \rightleftharpoons \mathrm{HF(s)}$
          \end{itemize}
          where $h^{+}$ denotes the oxidising half of the salt redox couple. In
          practice the fluorine potential is set by a buffer (Be metal, or an
          imposed $\mathrm{HF/H_2}$ ratio in the cover gas) and channel F is
          written with an effective activity $a_\mathrm{F}$ lumping the fluoride
          activity and the redox potential:
\end{itemize}

\begin{equation}
    w_{\mathrm{F}} = k_{\mathrm{f}}^{+}\, a_{\mathrm{F}}\,
                     \cm{\mathrm{H}}\big|_\Gamma
                   - k_{\mathrm{f}}^{-}\,\cs{\mathrm{HF}}\big|_\Gamma
    \label{eq:model3_rate}
\end{equation}

\begin{itemize}
    \item The atomic flux out of the metal is now shared:
\end{itemize}

\begin{equation}
    -D^{\mathrm{m}}\nabla \cm{\mathrm{H}}\cdot\mathbf{n}\big|_\Gamma
    = 2\,w_{\mathrm{rec}} + w_{\mathrm{F}}
    \label{eq:model3_total_flux}
\end{equation}

\begin{itemize}
    \item \KEY{Channel R is second order in $\cm{\mathrm{H}}$, channel F is
          first order.} The carriers they feed obey different transport and
          different release physics at the free surface. Define the branching
          ratio:
\end{itemize}

\begin{equation}
    \Bra
    \equiv \frac{w_{\mathrm{F}}}{2\,w_{\mathrm{rec}}}
    \;\xrightarrow[\text{far from equilibrium}]{}\;
    \frac{k_{\mathrm{f}}^{+} a_{\mathrm{F}}}
         {2\,k_{\mathrm{r}}^{+}\,\cm{\mathrm{H}}\big|_\Gamma}
    \label{eq:branching}
\end{equation}

\begin{itemize}
    \item Consequences to emphasise --- this is the core argument of the paper:
          \begin{itemize}
              \item $\Bra$ depends \emph{inversely} on the interfacial loading.
                    At low loading the fluorination channel dominates and the
                    salt-side inventory scales linearly with $\cm{\mathrm{H}}$;
                    at high loading recombination takes over and the scaling
                    becomes quadratic. Define the apparent interfacial exponent
                    \begin{equation*}
                        n \equiv \frac{\partial \ln \left(\text{total salt-side H}\right)}
                                      {\partial \ln \cm{\mathrm{H}}},
                        \qquad 1 \le n \le 2 .
                    \end{equation*}
                    LTE requires $n$ to be a constant chosen in advance ($n=2$
                    for Sieverts/Henry, $n=1$ for a Henry/Henry-like closure).
                    Here $n$ is a \emph{solution-dependent} quantity that drifts
                    during a transient. No fixed algebraic interface law
                    reproduces this.
              \item $\Bra$ depends on the salt redox state through
                    $a_\mathrm{F}$, which is itself dynamic: HF production,
                    corrosion of the container, and buffer depletion all move
                    it. Interface behaviour is therefore coupled to salt
                    chemistry, not a material property of the pair.
              \item Even when both channels are individually fast
                    ($\Da \gg 1$ for each), the \emph{partition} between them is
                    set by the ratio of rate constants, not by thermodynamics.
                    Fast kinetics does not rescue LTE here; it only makes each
                    channel individually equilibrated while the split between
                    them remains kinetic.
          \end{itemize}

    \item \textbf{Honest caveat} (write this --- it pre-empts the obvious
          referee objection and makes the argument stronger, not weaker): one
          \emph{could} close the interface by imposing full local chemical
          equilibrium among all interfacial species, a Gibbs minimisation at
          $\Gamma$ given the fluorine potential. That closure is
          \begin{enumerate}[label=(\roman*)]
              \item not what any macroscopic hydrogen transport code implements;
              \item valid only if all interconversions are fast relative to
                    transport, which is exactly what is in doubt
                    \cite{hodille_molecular_2022,calderoni_measurement_2008};
              \item still silent on how a given net flux splits between
                    carriers, since equilibrium fixes concentrations and not
                    fluxes.
          \end{enumerate}
\end{itemize}

% ---------------------------------------------------------------
\subsection{Multiple isotopes: isotopologue channels}
\label{sec:isotopes}
% ---------------------------------------------------------------

\begin{itemize}
    \item Extend to $N_\mathrm{iso}$ mobile atomic species in the metal.
          Recombination now populates every isotopologue; fluorination every
          fluoride. For $\{\mathrm{H,T}\}$:
\end{itemize}

\begin{align}
    w_{\mathrm{HH}} &= k_{\mathrm{HH}}^{+}\left(\cm{\mathrm{H}}\right)^{2}
                     - k_{\mathrm{HH}}^{-}\,\cs{\mathrm{H_2}}
    \label{eq:iso_HH}\\
    w_{\mathrm{HT}} &= 2\,k_{\mathrm{HT}}^{+}\,\cm{\mathrm{H}}\,\cm{\mathrm{T}}
                     - k_{\mathrm{HT}}^{-}\,\cs{\mathrm{HT}}
    \label{eq:iso_HT}\\
    w_{\mathrm{TT}} &= k_{\mathrm{TT}}^{+}\left(\cm{\mathrm{T}}\right)^{2}
                     - k_{\mathrm{TT}}^{-}\,\cs{\mathrm{T_2}}
    \label{eq:iso_TT}
\end{align}

\begin{align}
    w_{\mathrm{HF}} &= k_{\mathrm{f,H}}^{+} a_{\mathrm{F}} \cm{\mathrm{H}}
                     - k_{\mathrm{f,H}}^{-}\,\cs{\mathrm{HF}}
    \label{eq:iso_HF}\\
    w_{\mathrm{TF}} &= k_{\mathrm{f,T}}^{+} a_{\mathrm{F}} \cm{\mathrm{T}}
                     - k_{\mathrm{f,T}}^{-}\,\cs{\mathrm{TF}}
    \label{eq:iso_TF}
\end{align}

\begin{itemize}
    \item The factor 2 in Eq.~(\ref{eq:iso_HT}) is the statistical degeneracy of
          the mixed pair. \TODO{state the convention explicitly.}

    \item \KEY{The counting argument.} Two mobile species on the metal side feed
          five carriers on the salt side:
          \begin{equation*}
              N_\mathrm{iso} \;\longrightarrow\;
              \underbrace{\tfrac{1}{2}N_\mathrm{iso}(N_\mathrm{iso}+1)}_{\text{molecular}}
              + \underbrace{N_\mathrm{iso}}_{\text{fluoride}},
              \qquad
              \mathrm{(H,T)}:\; 2 \rightarrow 3 + 2 = 5 .
          \end{equation*}
          LTE supplies one scalar constraint per transported species pair. With
          two metal-side and five salt-side interfacial unknowns the per-species
          Dirichlet closure is underdetermined. Worse, it is not merely
          underdetermined but structurally wrong in form: by
          Eq.~(\ref{eq:iso_HT}) the HT concentration is set by the
          \emph{product} $\cm{\mathrm{H}}\cm{\mathrm{T}}$, so the
          metal~$\rightarrow$~salt map is bilinear and non-diagonal. A condition
          of the form $c_i/K_{S,i} = f(c_i')$ cannot represent an off-diagonal
          coupling. LTE is not inaccurate here; it is ill-posed.

    \item Isotope exchange,
          $\mathrm{H_2} + \mathrm{T_2} \rightleftharpoons 2\,\mathrm{HT}$, is
          slow homogeneously but is catalysed by the metal surface. In this
          framework it emerges automatically from
          Eqs.~(\ref{eq:iso_HH})--(\ref{eq:iso_TT}): a $\mathrm{T_2}$ molecule
          may dissociate at $\Gamma$, deposit T into the lattice, and the lattice
          T may leave paired with an H. Interfacial scrambling is thus a
          prediction of the model, not an input.

    \item Two observable consequences, both invisible to LTE:
          \begin{itemize}
              \item \textbf{Isotope swamping.} The T flux depends on the H
                    inventory: raising $\cm{\mathrm{H}}$ shifts T from the
                    $\mathrm{T_2}$ channel to the statistically favoured HT
                    channel and changes the total T throughput. LTE, being
                    diagonal, predicts no such coupling.
              \item \textbf{Interfacial fractionation.} Because the R and F
                    channels have different isotope dependences
                    ($k_\mathrm{HH}$ vs $k_\mathrm{HT}$ vs $k_\mathrm{TT}$;
                    $k_\mathrm{f,H}$ vs $k_\mathrm{f,T}$), the isotopic
                    composition of the permeating flux differs from that of the
                    metal, and differs between the molecular and the fluoride
                    carrier. Define
                    $\alpha = (\mathrm{T/H})_\mathrm{flux} /
                              (\mathrm{T/H})_\mathrm{metal}$
                    and report it as a model output.
          \end{itemize}

    \item Detailed balance for the isotopologue channels: the ratios
          $k^{+}/k^{-}$ are constrained by the isotopologue equilibrium
          constants, which are close to but not exactly the classical
          statistical values because of zero-point energy. Collected in
          App.~\ref{app:balance}.
\end{itemize}

% ---------------------------------------------------------------
\subsection{Dimensionless criteria for LTE validity}
\label{sec:damkohler}
% ---------------------------------------------------------------

\begin{itemize}
    \item \textbf{Within a channel}: compare interfacial equilibration with bulk
          transport.
\end{itemize}

\begin{equation}
    \Da_r = \frac{k_r^{+}\,L}{D}
    \label{eq:damkohler}
\end{equation}

\begin{itemize}
    \item $\Da_r \gg 1$ implies channel $r$ is locally equilibrated and its
          algebraic LTE form is recovered (Secs.~\ref{sec:model1},
          \ref{sec:model2}). For asymmetric systems evaluate one $\Da$ per side
          ($D_A, L_A$ vs $D_B, L_B$); the smaller governs.

    \item Caveat carried over from \cite{silva-solis_hydrogen_2026}:
          $\Da \gg 1$ is \emph{necessary but not sufficient}. Under a net flux
          the steady state departs from equilibrium by a term proportional to
          $\phi/k^{+}$, which LTE sets to zero by construction.

    \item \textbf{Between channels}: the branching ratio $\Bra$ of
          Eq.~(\ref{eq:branching}).
          \begin{itemize}
              \item $\Bra \rightarrow 0$ or $\Bra \rightarrow \infty$: one
                    channel dominates, an LTE form exists (Sieverts/Henry or
                    linear respectively).
              \item $\Bra = \mathcal{O}(1)$: no single algebraic law applies at
                    all, regardless of how large the $\Da$ are.
          \end{itemize}
          \TODO{present as a two-axis regime map $(\Da, \Bra)$. This figure is
          probably the single most useful deliverable of the paper.}

    \item \textbf{Across isotopes}: LTE requires the metal~$\rightarrow$~salt map
          to be diagonal. Quantify the departure by the fraction of the T flux
          carried by the mixed isotopologue,
          \begin{equation*}
              \chi_\mathrm{HT} = \frac{w_\mathrm{HT}}{w_\mathrm{HT} + 2 w_\mathrm{TT}} .
          \end{equation*}
          $\chi_\mathrm{HT} \rightarrow 0$ recovers a diagonal (per-isotope)
          description; $\chi_\mathrm{HT} = \mathcal{O}(1)$ forbids it. In a
          T-lean, H-rich system $\chi_\mathrm{HT} \rightarrow 1$.
\end{itemize}

% ===============================================================
\section{Implementation in \FESTIM}
\label{sec:implementation}
% ===============================================================

% ---------------------------------------------------------------
\subsection{Weak formulation}
\label{sec:weak}
% ---------------------------------------------------------------

\begin{itemize}
    \item Bulk equations per species per subdomain, multiplied by test functions
          $v \in H^1(\Omega)$ and integrated by parts. Interface terms appear
          naturally as surface integrals on $\Gamma$:
          \begin{equation*}
              \int_\Gamma \Big(\sum_r \nu_{ir} w_r\Big) v^{\mathrm{m}}_i \,\mathrm{d}S,
              \qquad
              \int_\Gamma \Big(\sum_r \mu_{\alpha r} w_r\Big) v^{\mathrm{s}}_\alpha \,\mathrm{d}S,
          \end{equation*}
          with $w_r$ given by Eq.~(\ref{eq:mass_action}).

    \item \KEY{No interfacial degrees of freedom are introduced}: all channels
          are algebraic in the trace values. This is the key simplification
          relative to the companion codimension-1 work
          \cite{festim_codim1_inprep}, and worth pointing out --- it means the
          present models require only submesh coupling through entity maps, not
          an additional function space on $\Gamma$.

    \item Implementation on submeshes in DOLFINx with
          \texttt{MixedFunctionSpace} and entity maps, following the
          multi-species architecture of \FESTIM{} v2 \cite{dark_festim_2026}.

    \item Nonlinearity: quadratic (recombination) and bilinear (isotopologue)
          terms; Newton with analytic Jacobian. Jacobian blocks in
          App.~\ref{app:weak}.
\end{itemize}

% ---------------------------------------------------------------
\subsection{Numerical considerations}
\label{sec:numerics}
% ---------------------------------------------------------------

\begin{itemize}
    \item Stiffness at large $\Da$: the interface conditions approach an
          algebraic constraint. \TODO{discuss conditioning and any scaling of the
          interface residual used to keep Newton well behaved.}
    \item Positivity: concentrations must stay non-negative for the quadratic
          terms to be meaningful. \TODO{note the safeguard used, if any.}
    \item Automatic recovery of LTE: no special-casing needed; the LTE solution
          is approached continuously as $k \rightarrow \infty$. Demonstrated
          numerically in Sec.~\ref{sec:verification}.
    \item Cost: report the overhead of $N_\mathrm{channel}$ surface assemblies
          relative to the LTE Dirichlet interface on the same mesh.
\end{itemize}

% ---------------------------------------------------------------
\subsection{Parameterisation}
\label{sec:params}
% ---------------------------------------------------------------

\begin{itemize}
    \item Detailed balance constrains the \emph{ratio} $k^{+}/k^{-}$ from
          thermodynamics; the \emph{magnitude} must come from elsewhere:
          \begin{enumerate}[label=(\roman*)]
              \item atomistics (DFT / MD) where available;
              \item fitting to permeation breakthrough transients;
              \item analogy with gas/metal recombination coefficients as a first
                    estimate for the R channel.
          \end{enumerate}

    \item For the F channel, $a_\mathrm{F}$ is tied to the salt redox state.
          \TODO{explain how it is set here --- imposed $\mathrm{HF/H_2}$ ratio in
          the cover gas, or a Be-metal buffer --- and state clearly which
          \HYPERION{} campaigns fix it.}

    \item Honest statement of uncertainty: FLiBe transport data for
          $\mathrm{H_2}$/HF and their isotopologues span orders of magnitude
          \cite{libra_2022}, which is precisely why an interface model whose
          parameters can be \emph{inferred} from permeation data is useful.

    \item \textbf{Important caveat}, developed in Sec.~\ref{sec:units}: part of
          that spread is not experimental scatter but a units/law inconsistency.
          A tabulated solubility carries $\mathrm{mol\,m^{-3}\,Pa^{-1}}$ or
          $\mathrm{mol\,m^{-3}\,Pa^{-1/2}}$ depending on whether Henry or
          Sieverts was assumed, and the two are not interconvertible by a fixed
          factor. Literature values therefore cannot be ingested without knowing
          the assumed carrier. Whenever a solubility is quoted in this paper, the
          assumed law is quoted with it.
\end{itemize}

% ===============================================================
\section{Verification}
\label{sec:verification}
% ===============================================================

% ---------------------------------------------------------------
\subsection{Recovery of LTE in the fast-kinetics limit}
\label{sec:lte_limit}
% ---------------------------------------------------------------

\begin{itemize}
    \item Model 1 on a 1D two-slab metal/metal problem: sweep $\Da$ over roughly
          six decades, plot the relative error in $c_A/c_B$ at steady state
          against the LTE (Sieverts/Sieverts) solution. Expect monotone
          first-order convergence in $1/\Da$; confirm the rate.
    \item Model 2 on a metal/liquid problem: same sweep, converging to the
          Sieverts/Henry solution, Eq.~(\ref{eq:lte_sh}). This is the more
          important of the two for this paper.
    \item Benchmark both against the existing \FESTIM{} LTE interface on the same
          mesh --- same answer to solver tolerance at large $\Da$.
\end{itemize}

% ---------------------------------------------------------------
\subsection{Analytical steady-state solutions}
\label{sec:analytical}
% ---------------------------------------------------------------

\begin{itemize}
    \item Model 1, 1D two-slab: closed form for $c_A|_\Gamma$, $c_B|_\Gamma$ and
          $\phi$ by combining flux continuity with Eq.~(\ref{eq:model1_flux}).
          The interface contributes an additive resistance $1/k^{+}$ in series
          with $L_A/D_A$ and $L_B/D_B$. Nice pedagogical framing: \emph{LTE is
          the zero-interfacial-resistance limit}.
    \item Model 3, steady state: solve the two-channel balance to get the
          apparent exponent $n$ as a function of $\Bra$ in closed form. This
          gives the analytical backbone of the regime-map figure.
    \item Full derivations in App.~\ref{app:analytical}.
\end{itemize}

% ---------------------------------------------------------------
\subsection{Conservation and consistency checks}
\label{sec:conservation}
% ---------------------------------------------------------------

\begin{itemize}
    \item Global atom balance: total H atoms (metal + salt + released) conserved
          to machine precision over a transient.
    \item Zero-flux test: with no driving gradient, the kinetic interface must
          relax to the LTE state. Confirms the detailed-balance constants are
          implemented consistently.
    \item MMS on the coupled two-domain problem with a manufactured interfacial
          source; verify spatial and temporal orders.
\end{itemize}

% ===============================================================
\section{Application: the \HYPERION{} experiment}
\label{sec:hyperion}
% ===============================================================

% ---------------------------------------------------------------
\subsection{System description and model setup}
\label{sec:hyperion_setup}
% ---------------------------------------------------------------

\begin{itemize}
    \item \HYPERION{} (HYdrogen PERmeatION) at the MIT PSFC
          \cite{hyperion_facility}: hydrogen isotopes permeate from a gas
          volume, through a Ni membrane, into a molten FLiBe pool held in a Ni
          container; the free surface of the salt is swept by a cover gas and
          the released flux is measured.
          \TODO{reproduce the geometry sketch; state dimensions, temperature
          range, sweep gas, and diagnostics (QMS / bubblers / ionisation
          chambers) --- confirm with the experimental team.}

    \item Three interfaces, each exercising a different part of the framework:
          \begin{enumerate}
              \item gas / Ni: existing dissociation--recombination BC;
              \item \textbf{Ni / FLiBe}: Sec.~\ref{sec:model3} --- competing R
                    and F channels. This is the interface of interest;
              \item FLiBe / cover gas: release of $\mathrm{H_2}$, HF and their
                    isotopologues from the free surface.
          \end{enumerate}

    \item Why this system is the right test: Calderoni et al.\
          \cite{calderoni_measurement_2008} measured T permeation through Ni into
          FLiBe and could not determine whether the tritium travelled as a
          fluoride-bound species, as HT, or as a mixture. A kinetic interface
          model turns that ambiguity into a parameter that can be inferred rather
          than assumed.

    \item Model reduction: justify the 1D treatment, or point to the
          multidimensional sidewall analysis \cite{hyperion_multidim} and state
          which correction is applied here.

    \item Baseline parameter table: $D$ and $K_S$ for H in Ni; $D$ and $K_H$ for
          $\mathrm{H_2}$ and HF in FLiBe; $k_\mathrm{r}$ and $k_\mathrm{f}$ at
          the Ni/FLiBe interface; $a_\mathrm{F}$. Give sources and uncertainty
          for each row and flag which are fitted.
          \TODO{add an ``assumed carrier'' column --- small editorial change that
          makes Sec.~\ref{sec:units} concrete.}
\end{itemize}

% ---------------------------------------------------------------
\subsection{Single isotope, single channel: LTE recovered}
\label{sec:hyperion_lte}
% ---------------------------------------------------------------

\begin{itemize}
    \item Run with the F channel switched off ($a_\mathrm{F} = 0$, strongly
          reducing salt). Sweep $k_\mathrm{r}$. Show:
          \begin{itemize}
              \item at large $\Da$ the kinetic model reproduces the \FESTIM{}
                    Sieverts/Henry LTE result for the permeation transient and
                    the steady flux;
              \item at moderate $\Da$ the breakthrough time lengthens and the
                    steady flux drops, by an amount set by the interfacial
                    resistance.
          \end{itemize}
    \item Establishes that the framework is a strict generalisation, and gives
          the reader a calibrated feel for how large $\Da$ must be.
\end{itemize}

% ---------------------------------------------------------------
\subsection{Competing channels: the apparent interfacial law}
\label{sec:hyperion_branching}
% ---------------------------------------------------------------

\begin{itemize}
    \item Turn on the F channel. Sweep $a_\mathrm{F}$ (redox state) and upstream
          pressure (loading) to move $\Bra$ across the regime map.

    \item \textbf{Headline figure}: apparent exponent $n$ against upstream
          pressure, for several $a_\mathrm{F}$. Show $n$ moving continuously
          between 1 and 2, i.e.\ the system is Henry-like at one end,
          Sieverts-like at the other, and neither in between. Overlay the two LTE
          predictions as horizontal lines to make the point visually.

    \item \textbf{Second figure}: flux speciation. Fraction of the released flux
          carried by $\mathrm{H_2}$ versus HF, against time and against
          $a_\mathrm{F}$. This is a \emph{direct output} of the model and is
          measurable (soluble/insoluble discrimination in bubbler counting),
          whereas LTE cannot even define it.

    \item \textbf{Third result}: fit $k_\mathrm{r}$, $k_\mathrm{f}$ to a
          \HYPERION{} permeation transient and report the inferred branching.
          Discuss identifiability honestly --- are $k_\mathrm{r}$ and
          $k_\mathrm{f}$ separately identifiable from a single transient, or only
          their ratio? Sensitivity study.

    \item \KEY{Strongest single piece of evidence in the paper:} show explicitly
          that no single choice of LTE law reproduces the measured transient
          across the redox sweep. Fit the Sieverts/Henry LTE model to the
          reducing case, then apply it unchanged to the oxidising case and
          quantify the error. Design the figure around this.
\end{itemize}

% ---------------------------------------------------------------
\subsection{H and T together: scrambling and fractionation}
\label{sec:hyperion_isotopes}
% ---------------------------------------------------------------

\begin{itemize}
    \item Five salt-side species ($\mathrm{H_2}$, HT, $\mathrm{T_2}$, HF, TF),
          two metal-side (H, T). Run a T-lean, H-rich case representative of a
          breeding blanket inventory, and a comparable-inventory case.

    \item \textbf{Result (a) --- isotope swamping.} Sweep the H partial pressure
          at fixed T source; show the T permeation flux and the T inventory in
          the salt changing. Report $\chi_\mathrm{HT}$. Contrast with the LTE
          prediction, which is by construction independent of the H inventory.

    \item \textbf{Result (b) --- interfacial fractionation.} Report
          $\alpha = (\mathrm{T/H})_\mathrm{flux}/(\mathrm{T/H})_\mathrm{metal}$
          for the molecular and the fluoride carrier separately. Show they
          differ, so the isotopic composition of the released flux depends on the
          redox state. Implication for tritium accountancy: a scheme calibrated
          on protium surrogates mis-attributes the T split.

    \item \textbf{Result (c) --- interfacial scrambling.} Start with pure
          $\mathrm{T_2}$ upstream and protium-loaded salt; show HT appearing in
          the released flux with no homogeneous exchange reaction in the model.
          The metal surface is the catalyst. Good place to note that this is
          qualitatively what the surrogate-species literature assumes but rarely
          models.

    \item \textbf{Result (d).} State plainly what an LTE run cannot produce here:
          there is no per-species Dirichlet condition to write for HT. Show the
          two ad hoc closures a modeller might reach for --- treat HT as an
          independent species with its own $K_H$, or lump all carriers into a
          single effective species --- and quantify how badly each does against
          the kinetic reference.
\end{itemize}

% ===============================================================
\section{Discussion}
\label{sec:discussion}
% ===============================================================

% ---------------------------------------------------------------
\subsection{Steady state versus thermodynamic equilibrium}
\label{sec:ss_vs_eq}
% ---------------------------------------------------------------

\begin{itemize}
    \item Restate the distinction of Sec.~\ref{sec:lte} now that the reader has
          seen it bite. Even at large $\Da$ the net flux forces a departure from
          equilibrium of order $\phi/k^{+}$.
    \item Quantify for \HYPERION{}: at experimental fluxes and plausible $k$, how
          large is the interfacial concentration jump that LTE discards?
\end{itemize}

% ---------------------------------------------------------------
\subsection{A solubility constant has no law-independent units}
\label{sec:units}
% ---------------------------------------------------------------

\begin{itemize}
    \item \TODO{pin the exact location --- which table or equation of
          \cite{calderoni_measurement_2008} carries the mismatched units, and
          whether the inconsistency is between the stated units and the fitted
          pressure exponent, or between the units and the permeability
          $\Phi = D K$ derived from them. Write this paragraph so a referee can
          check it in one minute; a vague accusation will be rejected.}

    \item The argument, stated as a general point rather than as a complaint
          about one paper:
          \begin{itemize}
              \item Henry: $c = K_H P
                    \Rightarrow [K_H] = \mathrm{mol\,m^{-3}\,Pa^{-1}}$
              \item Sieverts: $c = K_S\sqrt{P}
                    \Rightarrow [K_S] = \mathrm{mol\,m^{-3}\,Pa^{-1/2}}$
          \end{itemize}
          The \emph{dimensions} of the reported constant encode the assumed
          dissolution law. There is no conversion factor between the two:
          $K_H = K_S^2/c$ holds only along a single isobar, so any
          ``conversion'' silently fixes a reference pressure and is wrong
          everywhere else. A tabulated solubility is therefore not a material
          property that can be quoted independently of a speciation assumption
          --- it is a property of the (material $+$ assumed carrier) pair.

    \item Why this matters here specifically. Calderoni et al.\ report a
          square-root dependence of the permeating flux on pressure for FLiBe
          \cite{calderoni_measurement_2008}, whereas Fukada et al.\ report a
          linear dependence for FLiNaK, and the review literature states the
          expectation that dissolved diatomic $\mathrm{H_2}$ in a molten salt
          should obey Henry's law \cite{forsberg_flibe_2019}.
          \VERIFY{both against the primary sources before writing; the contrast
          is currently taken from a secondary summary \cite{mit_thesis_flibe} and
          needs the Fukada reference proper.}

    \item \KEY{A square-root dependence in an ionic melt is not a measurement
          artefact.} It is the signature of \emph{dissociative} dissolution ---
          i.e.\ of exactly the fluoride-bound / atomic carrier that the same
          authors could not rule out. In the language of Sec.~\ref{sec:model3}, a
          system running with $\Bra \gg 1$ (fluorination-dominated) presents an
          apparent exponent $n \rightarrow 1$, i.e.\ a Sieverts-like square-root
          response of salt inventory to upstream pressure, while a system with
          $\Bra \ll 1$ presents $n \rightarrow 2$ and looks Henry-like. The
          pressure exponent measured in a permeation experiment is therefore a
          \emph{readout of the branching ratio}, not a fixed property of the
          salt.

    \item This reframes an apparent inconsistency in the literature as a physical
          result. The FLiBe and FLiNaK datasets need not disagree: they may
          simply have been taken at different redox states, hence at different
          $\Bra$, hence with different apparent exponents. Present this as a
          falsifiable prediction --- \emph{a redox sweep at fixed temperature
          should move the measured exponent continuously between 1 and 2}. That
          is a concrete, cheap experiment for \HYPERION{} and probably the most
          useful single thing this paper can hand the experimentalists.

    \item Honest statement of the weakness this creates for \emph{us}:
          \begin{itemize}
              \item the most-cited Ni/FLiBe permeation dataset cannot be ingested
                    into a transport code without first assuming the speciation
                    that this paper argues is unknown --- circular;
              \item we therefore cannot use \cite{calderoni_measurement_2008} as
                    an independent parameterisation of the salt-side transport
                    properties; we use it only as qualitative evidence that the
                    interface is not a simple recombination boundary;
              \item our own $k_\mathrm{r}$, $k_\mathrm{f}$ are consequently
                    \emph{fitted} to \HYPERION{} transients rather than taken
                    from the literature, and the resulting identifiability
                    question is discussed in Sec.~\ref{sec:hyperion_branching};
              \item anywhere we do quote a literature solubility, state the
                    assumed law alongside it.
          \end{itemize}

    \item Broader recommendation for the community (one or two sentences, no
          more, or it reads as lecturing): report the fitted pressure exponent
          and the salt redox state alongside any tabulated solubility, so that
          the number remains usable if the speciation interpretation changes.
\end{itemize}

% ---------------------------------------------------------------
\subsection{When can LTE still be used?}
\label{sec:when_lte}
% ---------------------------------------------------------------

\begin{itemize}
    \item A practical decision procedure --- this is what practitioners will
          actually cite:
          \begin{enumerate}
              \item Is there a single exchange channel? If not, LTE is
                    unavailable, independent of $\Da$.
              \item Is the metal~$\rightarrow$~liquid species map diagonal
                    (single isotope)? If not, LTE is ill-posed.
              \item If both yes: check $\Da \gg 1$, and check
                    $\phi/(k^{+}c) \ll 1$.
          \end{enumerate}
    \item \KEY{Emphasise the ordering: the structural questions come first.} The
          literature has focused on step 3 and skipped 1 and 2.
\end{itemize}

% ---------------------------------------------------------------
\subsection{Implications for tritium accountancy and blanket design}
\label{sec:implications}
% ---------------------------------------------------------------

\begin{itemize}
    \item Speciation controls where tritium goes: HF/TF is highly soluble and
          corrosive and stays in the salt; $\mathrm{H_2}$/HT is mobile and
          permeates out. A model that cannot predict the split cannot predict
          either the inventory or the release.
    \item Redox control as a design variable: the branching ratio is something an
          operator can move. Frame this as an actionable result.
    \item Consequence for surrogate experiments: protium- or deuterium-only
          campaigns constrain the channels but not the isotopologue partitioning;
          fractionation must be modelled to extrapolate to T.
\end{itemize}

% ---------------------------------------------------------------
\subsection{Limitations and future work}
\label{sec:limitations}
% ---------------------------------------------------------------

\begin{itemize}
    \item Interfaces carry no stored inventory here. Interface trapping and the
          codimension-1 formulation needed for it are the companion paper
          \cite{festim_codim1_inprep}; the W/Cu
          \cite{silva-solis_hydrogen_2026} and Be/BeO
          \cite{hodille_molecular_2022} cases belong there.

    \item $a_\mathrm{F}$ treated as a prescribed parameter: no coupled salt redox
          / corrosion model. Natural next step is two-way coupling to a
          thermochemical solver.

    \item \textbf{Parameter provenance.} Because the available FLiBe solubility
          data are reported under an assumed dissolution law
          (Sec.~\ref{sec:units}), the salt-side transport parameters used here
          are fitted rather than independently measured. The framework is
          therefore demonstrated to be \emph{consistent with} the \HYPERION{}
          observations, not validated against an independent parameterisation.
          Say this plainly --- a referee will raise it otherwise, and the honest
          version is not damaging: the whole point is that an independent
          parameterisation does not currently exist.

    \item Interface assumed planar and uniform; no natural convection in the
          salt; no wetting or gas-bubble effects at the free surface.

    \item Rate constants assumed Arrhenius and concentration independent.

    \item The framework is not specific to FLiBe or to fusion: any interface
          where the transported element changes chemical identity on crossing
          (LiPb, oxide layers, permeation barriers, electrolysis, catalysis) has
          the same structure.
\end{itemize}

% ===============================================================
\section{Conclusions}
\label{sec:conclusions}
% ===============================================================

\begin{itemize}
    \item Formulated a general kinetic interface framework for hydrogen isotope
          transport: reversible reaction channels on the interface obeying mass
          action, with detailed balance fixing rate ratios from thermodynamics.
          Implemented in \FESTIM{}.

    \item LTE is recovered analytically as the fast-kinetics limit of a
          \emph{single} channel; both the Sieverts/Sieverts and Sieverts/Henry
          interface laws in current use are limits of different channels of the
          same framework.

    \item Two dimensionless groups delimit validity: a Damk\"ohler number within
          a channel, and a branching ratio between channels. Large Damk\"ohler is
          necessary but not sufficient.

    \item For the \HYPERION{} Ni/FLiBe system, hydrogen leaving the metal
          partitions kinetically between molecular and fluoride carriers. The
          apparent interfacial law is neither Sieverts nor Henry but drifts
          between them with loading and salt redox state, so no fixed LTE closure
          reproduces the measured behaviour across a redox sweep.

    \item With H and T present, two metal-side species feed five salt-side
          carriers through a bilinear, non-diagonal map. A per-species LTE
          condition is not merely inaccurate but ill-posed. Isotope swamping,
          interfacial fractionation and surface-catalysed scrambling are
          predicted and are measurable.

    \item Flux speciation becomes a direct simulation output, enabling comparison
          with bubbler and mass-spectrometry diagnostics and turning a
          long-standing ambiguity about the Ni/FLiBe interface
          \cite{calderoni_measurement_2008} into an inferable parameter.

    \item Corollary for the property database: the pressure exponent measured in
          a salt permeation experiment is a readout of the interfacial branching
          ratio rather than a fixed property of the salt. This offers a physical
          reconciliation of the conflicting square-root and linear pressure
          dependences reported for fluoride melts, and yields a falsifiable
          prediction --- a redox sweep at fixed temperature should move the
          measured exponent continuously between 1 and 2.
\end{itemize}

% ===============================================================
\section*{Acknowledgements}
% ===============================================================

\TODO{acknowledgements.}

% ===============================================================
\appendix
% ===============================================================

% ---------------------------------------------------------------
\section{Detailed balance constraints}
\label{app:balance}
% ---------------------------------------------------------------

\begin{itemize}
    \item Model 1: derivation of $k^{+}/k^{-} = K_{S,B}/K_{S,A}$.
    \item Model 2: derivation of
          $k_\mathrm{r}^{+}/k_\mathrm{r}^{-} = K_H/K_S^2$, and the resulting
          equivalence with the Sieverts/Henry LTE condition.
    \item Model 3: constraint on the F channel from the fluorine potential /
          redox buffer; relation between $a_\mathrm{F}$ and the imposed
          $\mathrm{HF/H_2}$ ratio.
    \item Isotopologue channels: equilibrium constants for
          $\mathrm{H_2} + \mathrm{T_2} \rightleftharpoons 2\,\mathrm{HT}$,
          including the zero-point-energy departure from the classical
          statistical value $K = 4$.
\end{itemize}

% ---------------------------------------------------------------
\section{The LTE limit}
\label{app:limits}
% ---------------------------------------------------------------

\begin{itemize}
    \item Formal asymptotics for $k^{+} \rightarrow \infty$ at fixed ratio, for
          Models 1 and 2; leading-order correction in $1/\Da$.
    \item Demonstration that no such limit exists for Model 3 with
          $\Bra = \mathcal{O}(1)$.
\end{itemize}

% ---------------------------------------------------------------
\section{Analytical steady-state solutions}
\label{app:analytical}
% ---------------------------------------------------------------

\begin{itemize}
    \item Model 1, 1D two-slab: interfacial resistance in series with the bulk.
    \item Model 3, steady state: apparent exponent $n$ as a function of $\Bra$.
\end{itemize}

% ---------------------------------------------------------------
\section{Weak form derivations}
\label{app:weak}
% ---------------------------------------------------------------

\begin{itemize}
    \item Weak form for each model.
    \item Jacobian blocks for the Newton solver, including the bilinear
          isotopologue coupling terms.
    \item Notes on the DOLFINx \texttt{MixedFunctionSpace} / entity map
          implementation.
\end{itemize}

% ===============================================================
\bibliographystyle{elsarticle-num}
\bibliography{references}
% ===============================================================

\end{document}